% =============================================================
% Введение (2 страницы)
% =============================================================
\chapter*{Введение}
\phantomsection\addcontentsline{toc}{chapter}{Введение}

Подшипники скольжения являются одними из~наиболее распространённых и~ответственных узлов трения в~современном машиностроении. Они применяются в~турбоагрегатах, центробежных насосах, компрессорах, буровом оборудовании и~многих других машинах, работающих в~условиях высоких нагрузок, скоростей и~температур. Надёжность и~ресурс подшипникового узла во~многом определяют работоспособность машины в~целом, а~потери на~трение составляют существенную долю общих энергетических потерь механизма~\cite{hori2002,szeri1998}. В~связи с~этим задача повышения трибологических характеристик подшипников скольжения~-- снижения потерь на~трение, увеличения нагрузочной способности, улучшения условий смазки~-- остаётся актуальной на~протяжении десятилетий.

Одним из~перспективных направлений решения этой задачи является нанесение на~рабочую поверхность деталей машин дискретного микрорельефа~-- регулярных рельефных элементов заданной геометрии и~расположения. В~отличие от~случайной шероховатости, элементы микрорельефа создаются целенаправленно методами лазерной обработки, фотолитографии, электроэрозионной обработки и~другими технологиями. Регулярные элементы~-- лунки, канавки, выступы~-- формируют в~смазочном слое дополнительные гидродинамические эффекты, позволяющие улучшить несущую способность подшипника, снизить коэффициент трения и~температуру в~зоне контакта~\cite{senatore2018,joshi2018}. Данное направление тесно связано с~концепцией зелёной трибологии, направленной на~снижение энергопотребления, уменьшение выбросов и~продление срока службы узлов трения~\cite{bashmur2024b,bashmur2022}.

\textbf{Целевая аудитория и~место в~учебном плане.} Настоящий учебник предназначен для~студентов направлений подготовки 15.03.02 и~15.04.02 <<Технологические машины и~оборудование>> (бакалавриат и~магистратура). Материал учебника может использоваться при~изучении следующих дисциплин:
\begin{itemize}
\item <<Трение, износ и~смазка в~машинах>>;
\item <<Технология повышения износостойкости объектов нефтегазового комплекса>>;
\item <<Спецглавы механики жидкости и~газа>>;
\item <<Повышение износостойкости деталей технологических машин и~оборудования>>;
\item <<Зеленые компетенции в~различных сферах жизни и~профессиональной деятельности>>.
\end{itemize}

\textbf{Цели и~задачи изучения.} В~результате изучения материала учебника обучающийся должен:
\begin{itemize}
\item \textit{знать} основные типы подшипников скольжения, режимы трения и~смазки, виды детерминированных микроструктур рабочей поверхности и~технологии их~формирования;
\item \textit{знать} математические постановки гидродинамических задач смазочного слоя на~основе уравнения Рейнольдса, методы учёта кавитации, термовязкостных эффектов и~шероховатости;
\item \textit{уметь} формулировать расчётную задачу для~подшипника скольжения с~микрорельефом поверхности, выбирать параметры микроструктур и~выполнять численный расчёт характеристик смазочного слоя;
\item \textit{уметь} анализировать и~интерпретировать результаты расчёта, проводить сравнительную оценку гладкого и~текстурированного вариантов подшипника;
\item \textit{владеть} навыками самостоятельного выполнения инженерного расчёта подшипника скольжения с~микрорельефом, оформления результатов в~виде пояснительной записки.
\end{itemize}

\textbf{Структура учебника.} Учебник состоит из~пяти глав, выстроенных по~принципу <<от~простого к~сложному>>.

В~первой главе изложены общие сведения о~подшипниках скольжения и~дискретном микрорельефе рабочей поверхности: классификация подшипников, принцип работы, режимы трения, типы и~параметры микроструктур, технологии их~формирования, связь с~концепцией зелёной трибологии.

Во~второй главе сформулированы математические постановки гидродинамических задач смазочного слоя для~радиальных и~упорных подшипников скольжения, введена параметрическая модель микрорельефа и~описана массосохраняющая модель кавитации (JFO).

В~третьей главе описаны численные методы решения: дискретизация уравнения Рейнольдса, итерационные решатели, GPU-реализация и~алгоритм решения задачи кавитации.

В~четвёртой главе рассмотрены дополнительные эффекты, влияющие на~работу смазочного слоя: нелинейность, термовязкостные эффекты, шероховатость поверхности и~динамика смазочного слоя.

Пятая глава посвящена курсовому проектированию. В~ней сформулированы цели и~задачи курсовой работы, приведены указания к~выполнению, варианты заданий и~три~подробных примера расчёта: радиальный подшипник центробежного насоса, радиальный подшипник шарошечного долота и~упорный подшипник скольжения.

\textbf{Особенности методики.} Каждая глава теоретической части (главы~1--4) завершается контрольными вопросами для~самопроверки, позволяющими обучающемуся оценить степень усвоения материала. Пятая глава содержит подробные примеры расчётов с~пошаговым описанием алгоритма, исходными данными и~анализом результатов, а~также варианты заданий для~самостоятельного выполнения курсовой работы. Такая структура обеспечивает последовательное формирование теоретических знаний и~практических навыков.

\textbf{Связь с~другими дисциплинами.} Изучение материала учебника опирается на~знания курсов высшей математики (дифференциальные уравнения, численные методы), механики жидкости и~газа, деталей машин, материаловедения. Полученные знания и~навыки применяются в~профессиональной деятельности при~проектировании и~модернизации узлов трения скольжения технологических машин и~оборудования.
